% ===================================================================
%  TABLEAU N° 12 — La gare. --- Les chemins de fer. --- Les navires.
%
%  Traduction, faite en 2026, en francais standard du Quebec, sur
%  l'ido de texto/io. Regles de la colonne : voir 10-tableau-01.tex.
%
%  LES SIX CHOIX PROPRES A CE TABLEAU :
%    (15) indicateur ............ horaire
%    (14) bibliothèque .......... kiosque à journaux
%    (3)  sac de nuit ........... sac de voyage
%    (66) cabinets d'aisances ... toilettes
%    (65) trottoirs ............. quais
%    (33) employés de l'octroi .. employés de l'accise
%
%  (53) ET (65) SONT LE MEME MOT EN IDO, comme (48) et (52) au tableau
%  6. Le fac-simile ecrit « quai » pour le premier et « trottoirs »
%  pour le second ; l'ido ecrit « trotuaro » deux fois, et il s'agit
%  du meme ouvrage — celui le long duquel le train s'arrete. On le
%  suit : quai les deux fois. « Trottoir » reste, mais pour ce qu'il
%  nomme au Quebec, le bord de la rue (tableau 11, renvoi 46).
%
%  (33) SUIT L'IDO CONTRE LE FAC-SIMILE. Le livret francais dit
%  « octroi », qui etait la taxe municipale francaise percue aux
%  portes des villes ; l'ido dit « acizo », l'accise, qui est l'impot
%  sur les marchandises et qui existe ici. Comme au tableau 2 pour le
%  jeu de la grenouille : on traduit l'ido.
% ===================================================================

%%K t12-apar-1 apar
\VUsaut{4.0mm}
\VUtitre{15.0pt}{60}{TABLEAU N\textsuperscript{o} 12}
\VUsaut{2.4mm}
\VUpk{11.4pt}{40}{La gare. --- Les chemins de fer. --- Les navires.}
\VUsaut{3.4mm}
\textsc{Note.} --- \textit{Les tableaux horaires employés dans chaque pays sont
différents; nous prenons donc pour exemple les heures du} \VUgras{tableau
français.}

%%K t12-01-1 p
1. --- Je viens de voyager en Allemagne. Avant le départ, j'avais acheté un
\VUgras{horaire} \textsuperscript{(15)} pour connaître l'heure de départ des
trains. Après avoir constaté que le train le plus commode partait de Paris à
neuf heures du soir et arrivait à Strasbourg à une heure treize, j'ai préparé
mon voyage et, le lendemain, je me suis fait conduire à la gare.

%%K t12-02-1 p
2. --- Quand je suis entré dans le grand hall des départs, derrière un vieux
\VUgras{prêtre} \textsuperscript{(2)} de campagne qui portait à la main son
\VUgras{sac de voyage} \textsuperscript{(3)}, un grand nombre de
\VUgras{voyageurs} \textsuperscript{(1)} étaient déjà arrivés.

%%K t12-02-2 p
Comme je voulais envoyer une \VUgras{dépêche (télégramme)}
\textsuperscript{(5)}, j'ai demandé à un \VUgras{contrôleur} \textsuperscript{(6)}
de m'indiquer \VUgras{le bureau du télégraphe} \textsuperscript{(4)}.

%%K t12-03-1 p
3. --- Avant de passer dans la \VUgras{salle d'attente} \textsuperscript{(7)},
je suis allé prendre un \VUgras{billet} \textsuperscript{(9)} aller-retour.

%%K t12-04-1 p
4. --- Je n'ai pas attendu longtemps au \VUgras{guichet} \textsuperscript{(8)},
parce qu'il y avait peu de monde. Un \VUgras{soldat} \textsuperscript{(10)} qui
partait en permission a eu l'obligeance de me céder son tour, de sorte que j'ai
eu le temps de demander quelques renseignements au \VUgras{chef de gare}
\textsuperscript{(13)}, qui causait avec le \VUgras{commissaire}
\textsuperscript{(12)} de surveillance et le \VUgras{gendarme}
\textsuperscript{(11)} de service.

%%K t12-tit-1 sub
\VUpk{10.2pt}{40}{L'enregistrement des bagages.}
\VUsaut{3.0mm}

%%K t12-05-1 p
5. --- Après avoir acheté un journal au \VUgras{kiosque à journaux}
\textsuperscript{(14)}, j'ai fait peser et enregistrer mes \VUgras{bagages}
\textsuperscript{(17)}, qu'un \VUgras{porteur} \textsuperscript{(16)} venait
d'amener sur un \VUgras{chariot} \textsuperscript{(19)}; \VUgras{l'employé}
\textsuperscript{(22)} préposé à la \VUgras{bascule} \textsuperscript{(21)} m'a
appris que ma malle pesait trente-cinq kilos, ce qui faisait un
\VUgras{excédent} de cinq kilos, pour lequel je devais payer un supplément au
\VUgras{guichet des bagages} \textsuperscript{(23)}.

%%K t12-06-1 p
6. --- Comme j'étais trop en avance, j'ai eu le loisir d'observer la circulation
des voyageurs dans la gare. Les uns déposaient ou reprenaient leurs menus
bagages à la \VUgras{consigne} \textsuperscript{(25)}; d'autres consultaient les
\VUgras{tableaux horaires} \textsuperscript{(26)}. Les voyageurs qui arrivaient
se pressaient vers la \VUgras{sortie} \textsuperscript{(32)}, leur
\VUgras{valise} \textsuperscript{(24)} à la main. Quelques-uns attendaient à la
\VUgras{distribution des bagages} \textsuperscript{(27)} et présentaient leur
\VUgras{bulletin} \textsuperscript{(29)} pour retirer leurs malles,
\VUgras{caisses} \textsuperscript{(28)} ou \VUgras{paniers}
\textsuperscript{(30)}.

%%K t12-07-1 p
7. --- \VUgras{Les employés de l'accise} \textsuperscript{(33)} inspectaient
soigneusement les colis pour vérifier s'il y avait quelque chose à déclarer.

%%K t12-08-1 p
8. --- Certains voyageurs se dirigeaient vers le \VUgras{buffet}
\textsuperscript{(34)}, où le \VUgras{garçon} \textsuperscript{(35)}
s'empressait de les servir. Ils devaient se dépêcher, parce que le
\VUgras{train} \textsuperscript{(36)}, qui est un express, ne s'arrête pas
longtemps en gare.

%%K t12-09-1 p
9. --- Je me suis rendu ensuite sur le quai de départ et j'ai cherché un
\VUgras{wagon} \textsuperscript{(37)} direct pour ne pas être obligé de changer
de voiture en cours de route. Je suis monté dans une voiture à couloir, qui
comprend un \VUgras{compartiment} \textsuperscript{(38)} pour fumeurs et un
compartiment réservé aux « dames seules ».

%%K t12-tit-2 sub
\VUpk{10.2pt}{40}{Le départ de nuit.}
\VUsaut{3.0mm}

%%K t12-10-1 p
10. --- Après avoir gravi le \VUgras{marchepied} \textsuperscript{(40)}, refermé
la \VUgras{portière} \textsuperscript{(39)}, relevé le \VUgras{store}
\textsuperscript{(41)} et placé mes menus bagages dans le \VUgras{filet}
\textsuperscript{(43)}, j'ai pris place sur la \VUgras{banquette}
\textsuperscript{(42)}. En voyant le \VUgras{lampiste} \textsuperscript{(44)}
allumer les lampes, j'ai pensé que je devrais passer la nuit dans le wagon et
j'ai appelé l'employé chargé de la location des \VUgras{oreillers}
\textsuperscript{(45)} pour lui en demander un.

%%K t12-11-1 p
11. --- Le contrôleur est venu vérifier les billets. Je lui ai demandé pourquoi
nous n'étions pas encore partis, puisque c'était l'heure. Il m'a répondu que le
\VUgras{chef de train} \textsuperscript{(46)} était occupé à faire placer des
vélos dans le \VUgras{fourgon des bagages} \textsuperscript{(47)}. De plus, on
n'avait pas encore changé toutes les bouillottes \textsuperscript{(48)}.

%%K t12-12-1 p
12. --- Mais nous allions bientôt partir, parce que le \VUgras{chauffeur}
\textsuperscript{(50)}, sur son \VUgras{tender} \textsuperscript{(49)}, prenait
du charbon pour activer le feu de la \VUgras{locomotive} \textsuperscript{(54)},
et \VUgras{le mécanicien} \textsuperscript{(51)} n'attendait que le signal du
\VUgras{sous-chef de gare} \textsuperscript{(52)}, qui était sur le
\VUgras{quai} \textsuperscript{(53)}.

%%K t12-13-1 p
13. --- La \VUgras{chaudière} \textsuperscript{(56)} de la locomotive grondait
sourdement; la \VUgras{cheminée} \textsuperscript{(55)} vomissait des torrents de
fumée mêlée de \VUgras{vapeur} \textsuperscript{(60)}. Un coup de
\VUgras{sifflet} \textsuperscript{(59)} strident a retenti et les
\VUgras{pistons} \textsuperscript{(58)} se sont mis en mouvement, actionnant les
\VUgras{roues} \textsuperscript{(57)} de la lourde machine.

%%K t12-tit-3 sub
\VUsaut{3.0mm}
\VUpk{10.2pt}{40}{En route!}
\VUsaut{3.0mm}

%%K t12-14-1 p
14. --- La vitesse du train, filant sur les \VUgras{rails} \textsuperscript{(64)},
s'est accélérée; les \VUgras{quais} \textsuperscript{(65)} ont disparu; nous
avons dépassé les \VUgras{toilettes} \textsuperscript{(66)}, puis des wagons
garés sur une autre \VUgras{voie} \textsuperscript{(63)} : des
\VUgras{wagons-lits} \textsuperscript{(67)}, un \VUgras{wagon-restaurant}
\textsuperscript{(68)}, un \VUgras{wagon-poste} \textsuperscript{(69)}.

%%K t12-noto-1 noto
\VUnotes{143.2mm}{%
\textit{(*) Note.} --- \textit{Il va de soi que la description du voyage suit
la frontière d'avant-guerre.}%
}

%%K t12-15-1 p
15. --- Puis c'est la \VUgras{gare des marchandises} \textsuperscript{(71)} qui
a passé devant nos yeux. Sous les hangars, des employés chargeaient les
\VUgras{marchandises} \textsuperscript{(72)} expédiées en petite vitesse. Des
\VUgras{hommes d'équipe} \textsuperscript{(76)}, près du \VUgras{château d'eau}
\textsuperscript{(73)}, manœuvraient des \VUgras{wagons à bestiaux}
\textsuperscript{(75)} et des \VUgras{wagons plateformes}
\textsuperscript{(79)} pour former un \VUgras{train de marchandises}
\textsuperscript{(80)}.

%%K t12-16-1 p
16. --- Nous avons franchi la dernière \VUgras{aiguille} \textsuperscript{(78)},
près de laquelle se tenait l'aiguilleur; nous avons dépassé le \VUgras{disque}
\textsuperscript{(86)} et la gare a disparu au loin.

%%K t12-17-1 p
17. --- Bientôt nous nous sommes engagés sur le \VUgras{pont de fer}
\textsuperscript{(74)} qui traverse le \VUgras{fleuve} \textsuperscript{(81)}.
Après ce pont, nous avons longé la rivière, sur laquelle de puissants
\VUgras{remorqueurs} \textsuperscript{(82)} traînaient de lourds
\VUgras{chalands} \textsuperscript{(83)}. Nous avons vu aussi un \VUgras{bateau
de voyageurs} \textsuperscript{(84)} au moment où il allait aborder au
\VUgras{ponton} \textsuperscript{(85)}.

%%K t12-18-1 p
18. --- Après avoir dépassé à toute vitesse plusieurs gares, nous avons traversé
quelques \VUgras{tunnels} \textsuperscript{(87)} et laissé derrière nous
d'innombrables \VUgras{maisonnettes} \textsuperscript{(88)} de
\VUgras{garde-barrières}. Bercé par le balancement du train, je me suis
profondément endormi.

%%K t12-tit-4 sub
\VUsaut{3.0mm}
\VUpk{10.2pt}{40}{La gare de Deutsch-Avricourt.}
\VUsaut{3.0mm}

%%K t12-19-1 p
19. --- J'ai été réveillé en sursaut par la voix d'un employé qui criait :
« Deutsch-Avricourt \textsuperscript{(*)}! Tout le monde descend! » Ce fut
l'inspection de la douane et toutes les fastidieuses formalités de la frontière.
J'ai dû régler ma montre sur l'\VUgras{horloge} \textsuperscript{(89)} de la
gare, qui avançait de cinquante-cinq minutes; à peine réveillé, je distinguais
difficilement la \VUgras{grande aiguille} \textsuperscript{(90)} de la
\VUgras{petite} \textsuperscript{(91)}; le \VUgras{cadran}
\textsuperscript{(92)} lui-même était tout brouillé par la brume du matin.

%%K t12-20-1 p
20. --- Pendant que les douaniers faisaient leur inspection, j'ai déchiffré en
bâillant les \VUgras{affiches} \textsuperscript{(93)} qui décorent les murs de la
gare. J'ai déjeuné pour passer le temps, j'ai consulté la carte du
\VUgras{réseau} \textsuperscript{(94)} allemand pour voir quelle distance me
séparait encore de Strasbourg, et je suis remonté dans mon wagon, réconforté par
l'espoir d'atteindre bientôt le but de mon voyage.
\VUsaut{6.0mm}
\VUfilet{22mm}
